\documentclass[10pt,a4paper]{amsart}
\usepackage{amsmath,amssymb,amsthm,graphicx}
% For commutative diagrams we use Xy-pic
\usepackage[all]{xy}

% We use 2cell for 2-commutative diagrams.
\xyoption{2cell}
% Theorem environments.
%
\theoremstyle{plain}
\newtheorem{theorem}[subsection]{Theorem}
\newtheorem{proposition}[subsection]{Proposition}
\newtheorem{lemma}[subsection]{Lemma}

\theoremstyle{definition}
\newtheorem{definition}[subsection]{Definition}
\newtheorem{example}[subsection]{Example}
\newtheorem{exercise}[subsection]{Exercise}
\newtheorem{situation}[subsection]{Situation}

\theoremstyle{remark}
\newtheorem{remark}[subsection]{Remark}
\newtheorem{remarks}[subsection]{Remarks}

\numberwithin{equation}{subsection}

% Macros
%
\def\lim{\mathop{\mathrm{lim}}\nolimits}
\def\colim{\mathop{\mathrm{colim}}\nolimits}
\def\Spec{\mathop{\mathrm{Spec}}}
\def\Hom{\mathop{\mathrm{Hom}}\nolimits}
\def\Ext{\mathop{\mathrm{Ext}}\nolimits}
\def\SheafHom{\mathop{\mathcal{H}\!\mathit{om}}\nolimits}
\def\SheafExt{\mathop{\mathcal{E}\!\mathit{xt}}\nolimits}
\def\Sch{\mathit{Sch}}
\def\Mor{\mathop{\mathrm{Mor}}\nolimits}
\def\Ob{\mathop{\mathrm{Ob}}\nolimits}
\def\Sh{\mathop{\mathit{Sh}}\nolimits}
\def\NL{\mathop{N\!L}\nolimits}
\def\CH{\mathop{\mathrm{CH}}\nolimits}
\def\proetale{{pro\text{-}\acute{e}tale}}
\def\etale{{\acute{e}tale}}
\def\QCoh{\mathit{QCoh}}
\def\Ker{\mathop{\mathrm{Ker}}}
\def\Im{\mathop{\mathrm{Im}}}
\def\Coker{\mathop{\mathrm{Coker}}}
\def\Coim{\mathop{\mathrm{Coim}}}

% Boxtimes
%
\DeclareMathSymbol{\boxtimes}{\mathbin}{AMSa}{"02}

%
% Macros for moduli stacks/spaces
%
\def\QCohstack{\mathcal{QC}\!\mathit{oh}}
\def\Cohstack{\mathcal{C}\!\mathit{oh}}
\def\Spacesstack{\mathcal{S}\!\mathit{paces}}
\def\Quotfunctor{\mathrm{Quot}}
\def\Hilbfunctor{\mathrm{Hilb}}
\def\Curvesstack{\mathcal{C}\!\mathit{urves}}
\def\Polarizedstack{\mathcal{P}\!\mathit{olarized}}
\def\Complexesstack{\mathcal{C}\!\mathit{omplexes}}
% \Pic is the operator that assigns to X its picard group, usage \Pic(X)
% \Picardstack_{X/B} denotes the Picard stack of X over B
% \Picardfunctor_{X/B} denotes the Picard functor of X over B
\def\Pic{\mathop{\mathrm{Pic}}\nolimits}
\def\Picardstack{\mathcal{P}\!\mathit{ic}}
\def\Picardfunctor{\mathrm{Pic}}
\def\Deformationcategory{\mathcal{D}\!\mathit{ef}}
\usepackage{cleveref}
\crefname{originaldefinition}{original definition}{original definitions}
\Crefname{originaldefinition}{Original definition}{Original definitions}
\makeatletter
\newlabel{situation-setup}{{02QL}{source}{}{sourcepointer-1}{}}
\newlabel{situation-setup@cref}{{[originaldefinition][0][]02QL}{[1][0][]source}{}{}{}}
\makeatother
\crefname{originallemma}{original lemma}{original lemmas}
\Crefname{originallemma}{Original lemma}{Original lemmas}
\makeatletter
\newlabel{section-splitting-principle}{{02UK}{source}{}{sourcepointer-2}{}}
\newlabel{section-splitting-principle@cref}{{[originallemma][0][]02UK}{[1][0][]source}{}{}{}}
\makeatother
\title{Computation of the alternating product of the Chern classes of the exterior powers of a locally free module}
\author{The Stacks Project Authors}
\date{stacks-project:chow.tex:local-snapshot-20261012; GNU FDL 1.2 or later}
\begin{document}
\maketitle
\noindent\textit{Editorial note.} All delivered mathematical text is the author's own, copied from the pinned Stacks Project chapter file chow.tex; no statement or proof step was rewritten, reordered or removed.\par\medskip
\noindent\textit{Editorial note.} Preamble provenance: the pinned chapter file opens with an input of the shared preamble and carries no preamble text of its own, so the project's single Stacks-wide preamble staged at sources/stacks-chow/preamble.tex is excerpted verbatim as two spans, lines 14-18 (Xy-pic setup) and lines 154-217 (the native theorem declarations and mathematical macros). Nothing was added, removed or redefined; the wrapper class and the amsmath/amssymb/amsthm/graphicx base packages are supplied by the assembler.\par\medskip
\noindent\textit{Editorial note.} The counted statement and proof contain no cross reference, no citation and no included file, so the delivered body is self contained.\par\medskip
\noindent\textit{Editorial note.} Canonical location of the counted proposition: Stacks tag 0FEE (sources/stacks-chow/pinned-original-source.tex line 12500).\par\medskip
\noindent\textit{Allowed standard prerequisites (editorial).}\begin{itemize}
\item The splitting principle for Chern classes of a locally free module and Chern roots. Same chapter, chow.tex, Section labelled section-splitting-principle (source line 8422), invoked at proof lines 12512-12514. Justifies replacing the computation on X by a computation in the polynomial ring on the Chern roots x\_1, ..., x\_r and supplies the formula c(wedge\textasciicircum{}n E) as a product over increasing multi-indices.
\item The base situation (S, delta) for bivariant operational Chow rings. Same chapter, Situation labelled situation-setup (source line 1372), invoked in the statement. Fixes the ring A\textasciicircum{}*(X) in which the identity is stated.
\end{itemize}\medskip
\ifdefined\hypertarget\hypertarget{sourcepointer-1}{}\fi
\noindent\textit{Source reference.} Original definition 02QL: sources/stacks-chow/pinned-original-source.tex, source lines 1372--1372. Same-chapter pointer into the author environment situation at source line 1372; invoked at proof line None of the counted proof. It is retained only as a navigation pointer; every load-bearing step is inside the delivered span.\par\medskip
\ifdefined\hypertarget\hypertarget{sourcepointer-2}{}\fi
\noindent\textit{Source reference.} Original lemma 02UK: sources/stacks-chow/pinned-original-source.tex, source lines 8422--8422. Same-chapter pointer into the author environment lemma at source line 8422; invoked at proof line 12514 of the counted proof. It is retained only as a navigation pointer; every load-bearing step is inside the delivered span.\par\medskip
\begin{lemma}
\label{lemma-compute-koszul}
Let $(S, \delta)$ be as in Situation \ref{situation-setup}. Let $X$
be a scheme locally of finite type over $S$. Let $\mathcal{E}$
be a locally free $\mathcal{O}_X$-module of rank $r$.
Then
$$
\prod\nolimits_{n = 0, \ldots, r} c(\wedge^n \mathcal{E})^{(-1)^n} =
1 - (r - 1)! c_r(\mathcal{E}) + \ldots
$$
\end{lemma}

\begin{proof}
By the splitting principle we can turn this into a calculation in the
polynomial ring on the Chern roots $x_1, \ldots, x_r$ of $\mathcal{E}$. See
Section \ref{section-splitting-principle}. Observe that
$$
c(\wedge^n \mathcal{E}) =
\prod\nolimits_{1 \leq i_1 < \ldots < i_n \leq r}
(1 + x_{i_1} + \ldots + x_{i_n})
$$
Thus the logarithm of the left hand side of the equation in the lemma is
$$
-
\sum\nolimits_{p \geq 1}
\sum\nolimits_{n = 0}^r
\sum\nolimits_{1 \leq i_1 < \ldots < i_n \leq r}
\frac{(-1)^{p + n}}{p}(x_{i_1} + \ldots + x_{i_n})^p
$$
Please notice the minus sign in front. However, we have
$$
\sum\nolimits_{p \geq 0}
\sum\nolimits_{n = 0}^r
\sum\nolimits_{1 \leq i_1 < \ldots < i_n \leq r}
\frac{(-1)^{p + n}}{p!}(x_{i_1} + \ldots + x_{i_n})^p
=
\prod (1 - e^{-x_i})
$$
Hence we see that the first nonzero term in our Chern class
is in degree $r$ and equal to the predicted value.
\end{proof}

\end{document}
